\honest{Einstein's Arrogance}{Eliezer Yudkowsky}{2007}

In 1919 Eddington went to observe an eclipse and test a prediction of General Relativity. A journalist, I say, asked Einstein what he would do if the observations failed to match his theory. Einstein replied: ``Then I would feel sorry for the good Lord. The theory is correct.'' \nb{In the account usually cited, Ilse Rosenthal-Schneider's memoir, the question came from its author, then a philosophy student, and after Lorentz had telegraphed Einstein a first report that Eddington had found the effect. Einstein was asked what he would have said had the result gone the other way.}

That sounds foolhardy. ``Who can know that the theory is correct, in advance of experimental test?'' Of course, Einstein turned out to be right, and I try not to criticize people when they are right.

My defense uses the arithmetic of my previous post. To pick the correct hypothesis out of 100 million, you need about 27 bits of evidence. The evidence is needed before any justifying, ``just to locate the hypothesis in theory-space,'' and hunches obey the same rule. \nb{Charles Peirce made this counting argument in 1903: out of ``trillions of trillions of hypotheses,'' the physicist hits on the right one within a dozen guesses, which chance cannot explain. He concluded that the guessing faculty is ``not strong enough to be oftener right than wrong.'' The post does not mention him.}

So when the equations first ``popped into his head,'' Einstein must already have held the evidence that singled them out. \nb{Einstein reached the final equations in 1915 after two years on an earlier version, the Entwurf theory, which failed to account for most of the anomaly in Mercury's orbit. In March 1914 he wrote of that theory that he no longer doubted ``the correctness of the whole system, regardless of whether the observation of the solar eclipse will be successful or not.'' It was wrong. The post's argument starts from the theory being right, so it cannot tell such confidence from justified confidence until the test is in.} I name none of the evidence he held. \nb{The best-known piece: in November 1915 the finished theory accounted for the whole anomaly in Mercury's orbit.}

How likely is it, I ask, that Einstein had exactly 29.5 bits, just enough to find a 29.3-bit theory and assign it 55\%? ``Not likely!'' So he ``probably had enough evidence to be damn sure.'' \nb{To be 99\% sure takes only about 6.6 bits more. The question is whether his evidence fell somewhere in that range, not at one point, and the post gives no reason to think that unlikely.} And since the brain is inefficient, he probably had far more evidence than ``a perfect Bayesian'' would need. \nb{That is a claim about what a perfect reasoner could conclude, not about how Einstein reached his own confidence.}

So the remark is not appalling. ``And remember that General Relativity was correct, from all that vast space of possibilities.''
